\section{上下端が変わる積分}\label{節:変数上下端}
$\int_a^x f(t)\dd t$ の微分が分かれば，上端や下端が別の関数になっている場合も合成関数の微分で扱える。$f$ がある区間で連続で，$u(x),v(x)$ が微分可能，その値がこの区間内にあるとする。区間内の固定した点 $c$ を選び，$A(s)=\int_c^s f(t)\dd t$ とおけば，
\[
 \int_{u(x)}^{v(x)}f(t)\dd t=A(v(x))-A(u(x)).
\]
両辺を微分し，$A'=f$ を使うと，
\[
 \boxed{\frac{\dd}{\dd x}\int_{u(x)}^{v(x)}f(t)\dd t
 =f(v(x))v'(x)-f(u(x))u'(x)}
\]
を得る。下端からの寄与に負号が付くのは，区間の加法性でその部分を引いているからである。$u(x)$ と $v(x)$ の大小が途中で入れ替わっても，向き付きの積分の約束によりこの式を使える。

\begin{example}{原始関数の式を求めずに微分する}
$G(x)=\int_x^{x^2}\e^{t^2}\dd t$ を微分せよ。
\end{example}
被積分関数 $\e^{t^2}$ はすべての実数で連続である。上端 $x^2$ の微分は $2x$，下端 $x$ の微分は $1$ だから，
\[
 G'(x)=\e^{(x^2)^2}\cdot2x-\e^{x^2}\cdot1
       =2x\e^{x^4}-\e^{x^2}.
\]
上端を代入してから内側の関数の微分を掛けるので，最初の指数は $x^4$ になる。$\e^{t^2}$ の原始関数を具体的な式で表す必要はない。

ただし，この公式では被積分関数を $f(t)$ とし，積分の外側の変数 $x$ がその中には現れないとしている。たとえば
\[
 H(x)=\int_0^x xt\dd t
      =x\int_0^x t\dd t=\frac{x^3}{2}
\]
では，$H'(x)=\frac{3x^2}{2}$ となる。上端を代入した $x^2$ だけでは足りない。積分の間は $x$ を定数として外に出せても，積分全体を $x$ で微分するときには，その因子も変わるからである。

